% Ad Familiares 7.16 --- headnote
% ad-familiares-07-16, November 54 BC, Rome

\textit{Cicero to C. Trebatius Testa, written from Rome in late
November 54 BC. Trebatius, the young jurist whom Cicero had pressed
on Caesar for service in Gaul, has by this point spent a long and
evidently unhappy season at the front: he missed Caesar's brief
expedition to Britain in the summer, and is now settled in winter
quarters at Samarobriva (Amiens), Caesar's Gallic headquarters. The
letter is one of a long teasing sequence in \textit{Fam.} 7 in which
Cicero ribs his protégé for grumbling, for staying clear of Britain,
and for failing to make his fortune as quickly as Balbus had promised
he would.}

\textit{The tone is the dry, lawyerly banter of two old Roman
friends: a tag from the old tragedy \textit{Equus Troianus} (``they
grow wise too late''), a Greek word (\textit{philotheoros}, ``fond of
sight-seeing'') dropped in to deflate the heroics of the British
campaign, a Stoic paradox bent into a joke about Trebatius's bank
balance, and a closing thrust that grants him primacy as a jurist
--- but only in Samarobriva. Cn. Octavius, who appears in section 2,
is a mutual acquaintance whose dinner invitations Cicero declines
with mock-hauteur.}
